\section{Inference settings and R code listings}
\label{app:implementation}
\label{app:code}

\subsection{Inference settings}

\begin{table}[htbp]
\centering
\small
\caption{Model-specific IF2 search settings.}
\begin{tabularx}{0.95\textwidth}{
  >{\raggedright\arraybackslash}p{0.20\textwidth}
  >{\raggedright\arraybackslash}p{0.37\textwidth}
  X}
\toprule
Model & Starting values & IF2 perturbation scale \\
\midrule
Gamma-noise
& \(B_0\in\{2,4,6\}\),
\(\sigB\in\{0.10,0.30,0.45\}\)
& \(B_0:\operatorname{ivp}(0.20)\);
\(\sigB:0.05\) \\
Constant-\(B\)
& \(\Bconst\in\{1,2,3,4,5,6\}\)
& \(\Bconst:0.05\) \\
\bottomrule
\end{tabularx}
\end{table}

Each start used 600 IF2 iterations, 5,000 particles per iteration, and
geometric cooling with half-life fraction 0.5.  Each fitted start was then
evaluated with five independent 50,000-particle filters.  Finite likelihood
estimates were combined with \texttt{pomp::logmeanexp}, and the start with the
largest summary was selected.  One final 50,000-particle filter retained
trajectory ancestry; \texttt{filter\_traj()} then extracted the sampled Gamma
path used in both the recovery metrics and trajectory figures.
Because the Gamma model used nine starts and the constant model six, the total
search budgets were unequal.

\subsection{Supplied R files}

The executable source is retained under
\path{code/experiment4_source/}.  The index below locates each workflow stage;
longer exact excerpts from the central scripts follow the table.

\begin{table}[htbp]
\centering
\small
\caption{Experiment~4 source-file index.}
\begin{tabularx}{\textwidth}{
  >{\ttfamily\raggedright\arraybackslash}p{0.42\textwidth}
  X}
\toprule
\normalfont File & Purpose \\
\midrule
experiment\_config.R
& Scientific parameters, simulation design, IF2 settings, and starting grids. \\
model\_components.R
& Data-generating, Gamma-noise, and constant-\(B\) POMP definitions. \\
01\_generate\_shared\_data\_task.R
& Generate and save one accepted data set for both fitting models. \\
02\_run\_gamma\_task.R
& Multi-start IF2, post-fit evaluation, and final filtering for Gamma-noise. \\
03\_run\_constant\_task.R
& Corresponding fitting workflow for constant-\(B\). \\
04\_combine\_results.R
& Validate and combine task-level model outputs. \\
05\_compare\_models.R
& Verify pairing, calculate recovery summaries, and generate comparison figures. \\
06\_make\_convergence\_diagnostics.R
& Create selected-task IF2 trace diagnostics and tail summaries. \\
07\_generate\_selected\_B\_trajectory.R
& Preserve the prespecified task-117 illustrative trajectory artifact. \\
\path{09_regenerate_sampled_B_trajectories.R}
& Reconstruct and validate the 200 sampled Gamma paths from saved best-fit parameters. \\
\bottomrule
\end{tabularx}
\end{table}

The following listings restore the implementation detail used to define the
models, run the Gamma-noise inference, and calculate the paired recovery
summaries.  The provenance limitation in
Appendix~\ref{app:reproducibility} still applies.

\clearpage
\subsection{Experiment configuration}

\lstinputlisting[
  style=Rstyle,
  firstline=1,
  lastline=52,
  firstnumber=1,
  caption={Complete preserved Experiment~4 configuration from
  \texttt{config/experiment\_config.R}, copied as
  \texttt{code/experiment4\_source/experiment\_config.R}.},
  label={lst:full-config}
]{code/experiment4_source/experiment_config.R}

\subsection{Data-generating and fitting POMPs}

\lstinputlisting[
  style=Rstyle,
  firstline=5,
  lastline=183,
  firstnumber=5,
  caption={Observation template, measurement model, prescribed piecewise
  data-generating POMP, latent Gamma-noise fitting POMP, and static
  constant-\(B\) fitting POMP from \texttt{code/model\_components.R}.},
  label={lst:model-builders}
]{code/experiment4_source/model_components.R}

\subsection{IF2 fitting and post-fit filtering}

\lstinputlisting[
  style=Rstyle,
  firstline=59,
  lastline=104,
  firstnumber=59,
  caption={Gamma-model parameter perturbations and the corresponding
  \texttt{mif2()} call from
  \texttt{code/02\_run\_gamma\_task.R}.},
  label={lst:if2-fit}
]{code/experiment4_source/02_run_gamma_task.R}

\lstinputlisting[
  style=Rstyle,
  firstline=167,
  lastline=230,
  firstnumber=167,
  caption={Independent post-fit particle-filter evaluations and
  likelihood-scale combination.},
  label={lst:postfit-evaluation}
]{code/experiment4_source/02_run_gamma_task.R}

\clearpage
\lstinputlisting[
  style=Rstyle,
  firstline=257,
  lastline=315,
  firstnumber=257,
  caption={Best-fit selection, final particle filter with retained ancestry,
  and sampled-trajectory extraction from
  \texttt{code/02\_run\_gamma\_task.R}.},
  label={lst:sampled-trajectory-code}
]{code/experiment4_source/02_run_gamma_task.R}

\subsection{Paired recovery calculations}

\lstinputlisting[
  style=Rstyle,
  firstline=50,
  lastline=117,
  firstnumber=50,
  caption={Shared-data validation and task-level recovery calculations from
  \texttt{code/05\_compare\_models.R}.},
  label={lst:paired-calculations}
]{code/experiment4_source/05_compare_models.R}

\lstinputlisting[
  style=Rstyle,
  firstline=167,
  lastline=240,
  firstnumber=167,
  caption={Paired comparison and across-task summaries from
  \texttt{code/05\_compare\_models.R}.  The implementation field
  \texttt{bias\_all} is called task-level mean error in this report.},
  label={lst:paired-summaries}
]{code/experiment4_source/05_compare_models.R}
